\documentclass[11pt]{article}
\usepackage[T1]{fontenc}
\usepackage{lmodern}
\usepackage[landscape,margin=0.7in]{geometry}
\usepackage{booktabs,graphicx,amsmath}
\usepackage[hidelinks]{hyperref}
\begin{document}
\title{Supplementary results: A Causal Evaluation of PCA-Guided Sensor-Rate Allocation under Hard Communication Budgets}
\author{Anonymous authors}
\date{}
\maketitle
\noindent All entries are mean $\pm$ sample standard deviation over five joint
experimental seeds (42, 123, 456, 789, 1024). These seeds vary acquisition,
detector randomness, and the Isolation Forest fitting subset. Full Data is a
full-rate reference and is listed only once per protocol. TEP point and event
metrics here use binary detection (any fault). Event boundaries never cross
simulation runs. Delay is in samples and conditional on detection; missed events
are reported separately. A dash denotes an undefined delay. Payload bytes
assume four bytes per retained value and exclude network overhead. RMSE and MAE
use train-standardized values. NRMSE divides channel RMSE by the pooled test
range, using one for a constant channel; those ranges are used only for reporting.
The main paper uses threshold quantile 0.99. Other predeclared quantiles are
sensitivity results and do not select the primary configuration.
\clearpage
\begin{table}[p]
\centering
\small
\caption{SMD: method retrained. Detection metrics; delay is conditional and measured in samples.}
\resizebox{\textwidth}{!}{%
\begin{tabular}{@{}llccccccccccc@{}}
\toprule
Dataset & Method & Budget & Precision & Recall & FPR & Event recall & Missed events & Delay (samples) & FA/1000 & Point F1 & Specificity & Events \\
\midrule
SMD-1-1 & PCA-Triage & 0.3 & 0.318 $\pm$ 0.009 & 0.526 $\pm$ 0.050 & 0.118 $\pm$ 0.009 & 0.675 $\pm$ 0.112 & 2.6 $\pm$ 0.9 & 37.114 $\pm$ 22.254 & 5.141 $\pm$ 0.413 & 0.396 $\pm$ 0.020 & 0.882 $\pm$ 0.009 & 8.0 $\pm$ 0.0 \\
SMD-1-1 & Variance & 0.3 & 0.326 $\pm$ 0.013 & 0.540 $\pm$ 0.026 & 0.117 $\pm$ 0.012 & 0.700 $\pm$ 0.068 & 2.4 $\pm$ 0.5 & 21.933 $\pm$ 9.999 & 5.780 $\pm$ 0.876 & 0.406 $\pm$ 0.008 & 0.883 $\pm$ 0.012 & 8.0 $\pm$ 0.0 \\
SMD-1-1 & Threshold & 0.3 & 0.321 $\pm$ 0.015 & 0.523 $\pm$ 0.023 & 0.116 $\pm$ 0.004 & 0.775 $\pm$ 0.105 & 1.8 $\pm$ 0.8 & 26.102 $\pm$ 13.551 & 4.797 $\pm$ 0.438 & 0.398 $\pm$ 0.018 & 0.884 $\pm$ 0.004 & 8.0 $\pm$ 0.0 \\
SMD-1-1 & Uniform & 0.3 & 0.319 $\pm$ 0.012 & 0.532 $\pm$ 0.032 & 0.119 $\pm$ 0.002 & 0.725 $\pm$ 0.163 & 2.2 $\pm$ 1.3 & 22.172 $\pm$ 15.011 & 4.010 $\pm$ 0.366 & 0.399 $\pm$ 0.018 & 0.881 $\pm$ 0.002 & 8.0 $\pm$ 0.0 \\
SMD-1-1 & PCA-Triage & 0.5 & 0.317 $\pm$ 0.013 & 0.507 $\pm$ 0.023 & 0.114 $\pm$ 0.003 & 0.700 $\pm$ 0.168 & 2.4 $\pm$ 1.3 & 26.070 $\pm$ 17.873 & 6.159 $\pm$ 0.543 & 0.390 $\pm$ 0.017 & 0.886 $\pm$ 0.003 & 8.0 $\pm$ 0.0 \\
SMD-1-1 & Variance & 0.5 & 0.329 $\pm$ 0.016 & 0.535 $\pm$ 0.033 & 0.114 $\pm$ 0.008 & 0.825 $\pm$ 0.143 & 1.4 $\pm$ 1.1 & 14.523 $\pm$ 15.970 & 6.988 $\pm$ 0.456 & 0.407 $\pm$ 0.019 & 0.886 $\pm$ 0.008 & 8.0 $\pm$ 0.0 \\
SMD-1-1 & Threshold & 0.5 & 0.320 $\pm$ 0.014 & 0.509 $\pm$ 0.017 & 0.113 $\pm$ 0.007 & 0.725 $\pm$ 0.163 & 2.2 $\pm$ 1.3 & 28.300 $\pm$ 17.448 & 6.342 $\pm$ 1.036 & 0.393 $\pm$ 0.013 & 0.887 $\pm$ 0.007 & 8.0 $\pm$ 0.0 \\
SMD-1-1 & Uniform & 0.5 & 0.319 $\pm$ 0.008 & 0.518 $\pm$ 0.009 & 0.116 $\pm$ 0.003 & 0.850 $\pm$ 0.137 & 1.2 $\pm$ 1.1 & 10.901 $\pm$ 7.975 & 5.337 $\pm$ 0.234 & 0.395 $\pm$ 0.008 & 0.884 $\pm$ 0.003 & 8.0 $\pm$ 0.0 \\
SMD-1-1 & PCA-Triage & 0.7 & 0.318 $\pm$ 0.019 & 0.525 $\pm$ 0.034 & 0.118 $\pm$ 0.007 & 0.750 $\pm$ 0.153 & 2.0 $\pm$ 1.2 & 18.590 $\pm$ 12.642 & 7.992 $\pm$ 0.832 & 0.396 $\pm$ 0.023 & 0.882 $\pm$ 0.007 & 8.0 $\pm$ 0.0 \\
SMD-1-1 & Variance & 0.7 & 0.320 $\pm$ 0.014 & 0.529 $\pm$ 0.027 & 0.118 $\pm$ 0.007 & 0.875 $\pm$ 0.177 & 1.0 $\pm$ 1.4 & 8.527 $\pm$ 5.246 & 8.680 $\pm$ 0.574 & 0.398 $\pm$ 0.015 & 0.882 $\pm$ 0.007 & 8.0 $\pm$ 0.0 \\
SMD-1-1 & Threshold & 0.7 & 0.328 $\pm$ 0.010 & 0.544 $\pm$ 0.010 & 0.116 $\pm$ 0.004 & 0.950 $\pm$ 0.112 & 0.4 $\pm$ 0.9 & 4.358 $\pm$ 0.828 & 8.554 $\pm$ 0.693 & 0.409 $\pm$ 0.010 & 0.884 $\pm$ 0.004 & 8.0 $\pm$ 0.0 \\
SMD-1-1 & Uniform & 0.7 & 0.320 $\pm$ 0.010 & 0.517 $\pm$ 0.022 & 0.115 $\pm$ 0.004 & 0.925 $\pm$ 0.068 & 0.6 $\pm$ 0.5 & 5.686 $\pm$ 2.465 & 7.051 $\pm$ 0.644 & 0.395 $\pm$ 0.013 & 0.885 $\pm$ 0.004 & 8.0 $\pm$ 0.0 \\
SMD-1-1 & Full Data & 1.0 & 0.330 $\pm$ 0.010 & 0.537 $\pm$ 0.011 & 0.114 $\pm$ 0.004 & 0.875 $\pm$ 0.125 & 1.0 $\pm$ 1.0 & 4.967 $\pm$ 0.709 & 8.968 $\pm$ 0.552 & 0.409 $\pm$ 0.010 & 0.886 $\pm$ 0.004 & 8.0 $\pm$ 0.0 \\
\addlinespace
\bottomrule
\end{tabular}%
}
\label{supp-det-smd-method_retrained}
\end{table}
\begin{table}[p]
\centering
\small
\caption{SMD: method retrained. Distortion and mask-counted payload. Four bytes per value are assumed. RMSE uses standardized values.}
\resizebox{\textwidth}{!}{%
\begin{tabular}{@{}llcccccccc@{}}
\toprule
Dataset & Method & Budget & MAE & RMSE & NRMSE & Values & Payload bytes & Realized fraction & Max excess values \\
\midrule
SMD-1-1 & PCA-Triage & 0.3 & 0.258 $\pm$ 0.027 & 14.902 $\pm$ 7.847 & 0.016 $\pm$ 0.000 & 324660.0 $\pm$ 0.0 & 1298640.0 $\pm$ 0.0 & 0.300 $\pm$ 0.000 & 0.000 $\pm$ 0.000 \\
SMD-1-1 & Variance & 0.3 & 0.141 $\pm$ 0.002 & 2.375 $\pm$ 0.415 & 0.013 $\pm$ 0.000 & 324660.0 $\pm$ 0.0 & 1298640.0 $\pm$ 0.0 & 0.300 $\pm$ 0.000 & 0.000 $\pm$ 0.000 \\
SMD-1-1 & Threshold & 0.3 & 0.413 $\pm$ 0.028 & 71.459 $\pm$ 11.325 & 0.019 $\pm$ 0.000 & 324660.0 $\pm$ 0.0 & 1298640.0 $\pm$ 0.0 & 0.300 $\pm$ 0.000 & 0.000 $\pm$ 0.000 \\
SMD-1-1 & Uniform & 0.3 & 0.700 $\pm$ 0.014 & 116.545 $\pm$ 11.910 & 0.016 $\pm$ 0.000 & 324660.0 $\pm$ 0.0 & 1298640.0 $\pm$ 0.0 & 0.300 $\pm$ 0.000 & 0.000 $\pm$ 0.000 \\
SMD-1-1 & PCA-Triage & 0.5 & 0.154 $\pm$ 0.010 & 9.024 $\pm$ 2.230 & 0.012 $\pm$ 0.000 & 541101.0 $\pm$ 0.0 & 2164404.0 $\pm$ 0.0 & 0.500 $\pm$ 0.000 & 0.000 $\pm$ 0.000 \\
SMD-1-1 & Variance & 0.5 & 0.058 $\pm$ 0.001 & 1.258 $\pm$ 0.246 & 0.008 $\pm$ 0.000 & 541101.0 $\pm$ 0.0 & 2164404.0 $\pm$ 0.0 & 0.500 $\pm$ 0.000 & 0.000 $\pm$ 0.000 \\
SMD-1-1 & Threshold & 0.5 & 0.109 $\pm$ 0.021 & 13.987 $\pm$ 17.107 & 0.013 $\pm$ 0.001 & 541101.0 $\pm$ 0.0 & 2164404.0 $\pm$ 0.0 & 0.500 $\pm$ 0.000 & 0.000 $\pm$ 0.000 \\
SMD-1-1 & Uniform & 0.5 & 0.421 $\pm$ 0.084 & 83.863 $\pm$ 19.805 & 0.012 $\pm$ 0.000 & 541101.0 $\pm$ 0.0 & 2164404.0 $\pm$ 0.0 & 0.500 $\pm$ 0.000 & 0.000 $\pm$ 0.000 \\
SMD-1-1 & PCA-Triage & 0.7 & 0.039 $\pm$ 0.002 & 1.902 $\pm$ 0.455 & 0.006 $\pm$ 0.000 & 757541.0 $\pm$ 0.0 & 3030164.0 $\pm$ 0.0 & 0.700 $\pm$ 0.000 & 0.000 $\pm$ 0.000 \\
SMD-1-1 & Variance & 0.7 & 0.015 $\pm$ 0.000 & 0.852 $\pm$ 0.065 & 0.004 $\pm$ 0.000 & 757541.0 $\pm$ 0.0 & 3030164.0 $\pm$ 0.0 & 0.700 $\pm$ 0.000 & 0.000 $\pm$ 0.000 \\
SMD-1-1 & Threshold & 0.7 & 0.021 $\pm$ 0.000 & 0.562 $\pm$ 0.055 & 0.004 $\pm$ 0.000 & 757541.0 $\pm$ 0.0 & 3030164.0 $\pm$ 0.0 & 0.700 $\pm$ 0.000 & 0.000 $\pm$ 0.000 \\
SMD-1-1 & Uniform & 0.7 & 0.219 $\pm$ 0.043 & 55.388 $\pm$ 21.253 & 0.009 $\pm$ 0.000 & 757541.0 $\pm$ 0.0 & 3030164.0 $\pm$ 0.0 & 0.700 $\pm$ 0.000 & 0.000 $\pm$ 0.000 \\
SMD-1-1 & Full Data & 1.0 & 0.000 $\pm$ 0.000 & 0.000 $\pm$ 0.000 & 0.000 $\pm$ 0.000 & 1082202.0 $\pm$ 0.0 & 4328808.0 $\pm$ 0.0 & 1.000 $\pm$ 0.000 & 0.000 $\pm$ 0.000 \\
\addlinespace
\bottomrule
\end{tabular}%
}
\label{supp-dist-smd-method_retrained}
\end{table}
\clearpage
\begin{table}[p]
\centering
\small
\caption{SMD: fixed full data. Detection metrics; delay is conditional and measured in samples.}
\resizebox{\textwidth}{!}{%
\begin{tabular}{@{}llccccccccccc@{}}
\toprule
Dataset & Method & Budget & Precision & Recall & FPR & Event recall & Missed events & Delay (samples) & FA/1000 & Point F1 & Specificity & Events \\
\midrule
SMD-1-1 & PCA-Triage & 0.3 & 0.320 $\pm$ 0.016 & 0.514 $\pm$ 0.019 & 0.114 $\pm$ 0.006 & 0.625 $\pm$ 0.000 & 3.0 $\pm$ 0.0 & 42.000 $\pm$ 9.445 & 4.860 $\pm$ 0.319 & 0.394 $\pm$ 0.016 & 0.886 $\pm$ 0.006 & 8.0 $\pm$ 0.0 \\
SMD-1-1 & Variance & 0.3 & 0.325 $\pm$ 0.014 & 0.525 $\pm$ 0.010 & 0.114 $\pm$ 0.006 & 0.650 $\pm$ 0.056 & 2.8 $\pm$ 0.4 & 25.573 $\pm$ 8.877 & 5.955 $\pm$ 0.560 & 0.401 $\pm$ 0.014 & 0.886 $\pm$ 0.006 & 8.0 $\pm$ 0.0 \\
SMD-1-1 & Threshold & 0.3 & 0.312 $\pm$ 0.011 & 0.517 $\pm$ 0.013 & 0.119 $\pm$ 0.006 & 0.750 $\pm$ 0.088 & 2.0 $\pm$ 0.7 & 29.100 $\pm$ 17.772 & 4.909 $\pm$ 0.248 & 0.389 $\pm$ 0.011 & 0.881 $\pm$ 0.006 & 8.0 $\pm$ 0.0 \\
SMD-1-1 & Uniform & 0.3 & 0.320 $\pm$ 0.015 & 0.516 $\pm$ 0.016 & 0.115 $\pm$ 0.005 & 0.675 $\pm$ 0.112 & 2.6 $\pm$ 0.9 & 24.406 $\pm$ 9.291 & 4.242 $\pm$ 0.430 & 0.395 $\pm$ 0.016 & 0.885 $\pm$ 0.005 & 8.0 $\pm$ 0.0 \\
SMD-1-1 & PCA-Triage & 0.5 & 0.325 $\pm$ 0.011 & 0.524 $\pm$ 0.015 & 0.114 $\pm$ 0.004 & 0.650 $\pm$ 0.056 & 2.8 $\pm$ 0.4 & 27.607 $\pm$ 18.407 & 6.096 $\pm$ 0.428 & 0.401 $\pm$ 0.012 & 0.886 $\pm$ 0.004 & 8.0 $\pm$ 0.0 \\
SMD-1-1 & Variance & 0.5 & 0.328 $\pm$ 0.012 & 0.532 $\pm$ 0.011 & 0.114 $\pm$ 0.005 & 0.900 $\pm$ 0.056 & 0.8 $\pm$ 0.4 & 10.450 $\pm$ 9.932 & 7.486 $\pm$ 0.331 & 0.406 $\pm$ 0.011 & 0.886 $\pm$ 0.005 & 8.0 $\pm$ 0.0 \\
SMD-1-1 & Threshold & 0.5 & 0.317 $\pm$ 0.013 & 0.530 $\pm$ 0.013 & 0.119 $\pm$ 0.005 & 0.825 $\pm$ 0.112 & 1.4 $\pm$ 0.9 & 11.692 $\pm$ 6.490 & 7.816 $\pm$ 0.391 & 0.396 $\pm$ 0.013 & 0.881 $\pm$ 0.005 & 8.0 $\pm$ 0.0 \\
SMD-1-1 & Uniform & 0.5 & 0.324 $\pm$ 0.011 & 0.523 $\pm$ 0.011 & 0.114 $\pm$ 0.005 & 0.775 $\pm$ 0.056 & 1.8 $\pm$ 0.4 & 9.567 $\pm$ 3.312 & 5.639 $\pm$ 0.304 & 0.400 $\pm$ 0.010 & 0.886 $\pm$ 0.005 & 8.0 $\pm$ 0.0 \\
SMD-1-1 & PCA-Triage & 0.7 & 0.328 $\pm$ 0.010 & 0.531 $\pm$ 0.011 & 0.114 $\pm$ 0.004 & 0.775 $\pm$ 0.137 & 1.8 $\pm$ 1.1 & 9.450 $\pm$ 6.326 & 7.606 $\pm$ 0.731 & 0.406 $\pm$ 0.010 & 0.886 $\pm$ 0.004 & 8.0 $\pm$ 0.0 \\
SMD-1-1 & Variance & 0.7 & 0.330 $\pm$ 0.011 & 0.535 $\pm$ 0.012 & 0.114 $\pm$ 0.004 & 0.875 $\pm$ 0.125 & 1.0 $\pm$ 1.0 & 5.005 $\pm$ 0.805 & 8.806 $\pm$ 0.712 & 0.408 $\pm$ 0.011 & 0.886 $\pm$ 0.004 & 8.0 $\pm$ 0.0 \\
SMD-1-1 & Threshold & 0.7 & 0.328 $\pm$ 0.011 & 0.535 $\pm$ 0.011 & 0.115 $\pm$ 0.005 & 0.825 $\pm$ 0.112 & 1.4 $\pm$ 0.9 & 5.718 $\pm$ 1.262 & 8.722 $\pm$ 0.672 & 0.407 $\pm$ 0.011 & 0.885 $\pm$ 0.005 & 8.0 $\pm$ 0.0 \\
SMD-1-1 & Uniform & 0.7 & 0.327 $\pm$ 0.009 & 0.530 $\pm$ 0.009 & 0.114 $\pm$ 0.004 & 0.875 $\pm$ 0.088 & 1.0 $\pm$ 0.7 & 5.632 $\pm$ 1.850 & 6.854 $\pm$ 0.580 & 0.404 $\pm$ 0.008 & 0.886 $\pm$ 0.004 & 8.0 $\pm$ 0.0 \\
SMD-1-1 & Full Data & 1.0 & 0.330 $\pm$ 0.010 & 0.537 $\pm$ 0.011 & 0.114 $\pm$ 0.004 & 0.875 $\pm$ 0.125 & 1.0 $\pm$ 1.0 & 4.967 $\pm$ 0.709 & 8.968 $\pm$ 0.552 & 0.409 $\pm$ 0.010 & 0.886 $\pm$ 0.004 & 8.0 $\pm$ 0.0 \\
\addlinespace
\bottomrule
\end{tabular}%
}
\label{supp-det-smd-fixed_full_data}
\end{table}
\begin{table}[p]
\centering
\small
\caption{SMD: fixed full data. Distortion and mask-counted payload. Four bytes per value are assumed. RMSE uses standardized values.}
\resizebox{\textwidth}{!}{%
\begin{tabular}{@{}llcccccccc@{}}
\toprule
Dataset & Method & Budget & MAE & RMSE & NRMSE & Values & Payload bytes & Realized fraction & Max excess values \\
\midrule
SMD-1-1 & PCA-Triage & 0.3 & 0.258 $\pm$ 0.027 & 14.902 $\pm$ 7.847 & 0.016 $\pm$ 0.000 & 324660.0 $\pm$ 0.0 & 1298640.0 $\pm$ 0.0 & 0.300 $\pm$ 0.000 & 0.000 $\pm$ 0.000 \\
SMD-1-1 & Variance & 0.3 & 0.141 $\pm$ 0.002 & 2.375 $\pm$ 0.415 & 0.013 $\pm$ 0.000 & 324660.0 $\pm$ 0.0 & 1298640.0 $\pm$ 0.0 & 0.300 $\pm$ 0.000 & 0.000 $\pm$ 0.000 \\
SMD-1-1 & Threshold & 0.3 & 0.413 $\pm$ 0.028 & 71.459 $\pm$ 11.325 & 0.019 $\pm$ 0.000 & 324660.0 $\pm$ 0.0 & 1298640.0 $\pm$ 0.0 & 0.300 $\pm$ 0.000 & 0.000 $\pm$ 0.000 \\
SMD-1-1 & Uniform & 0.3 & 0.700 $\pm$ 0.014 & 116.545 $\pm$ 11.910 & 0.016 $\pm$ 0.000 & 324660.0 $\pm$ 0.0 & 1298640.0 $\pm$ 0.0 & 0.300 $\pm$ 0.000 & 0.000 $\pm$ 0.000 \\
SMD-1-1 & PCA-Triage & 0.5 & 0.154 $\pm$ 0.010 & 9.024 $\pm$ 2.230 & 0.012 $\pm$ 0.000 & 541101.0 $\pm$ 0.0 & 2164404.0 $\pm$ 0.0 & 0.500 $\pm$ 0.000 & 0.000 $\pm$ 0.000 \\
SMD-1-1 & Variance & 0.5 & 0.058 $\pm$ 0.001 & 1.258 $\pm$ 0.246 & 0.008 $\pm$ 0.000 & 541101.0 $\pm$ 0.0 & 2164404.0 $\pm$ 0.0 & 0.500 $\pm$ 0.000 & 0.000 $\pm$ 0.000 \\
SMD-1-1 & Threshold & 0.5 & 0.109 $\pm$ 0.021 & 13.987 $\pm$ 17.107 & 0.013 $\pm$ 0.001 & 541101.0 $\pm$ 0.0 & 2164404.0 $\pm$ 0.0 & 0.500 $\pm$ 0.000 & 0.000 $\pm$ 0.000 \\
SMD-1-1 & Uniform & 0.5 & 0.421 $\pm$ 0.084 & 83.863 $\pm$ 19.805 & 0.012 $\pm$ 0.000 & 541101.0 $\pm$ 0.0 & 2164404.0 $\pm$ 0.0 & 0.500 $\pm$ 0.000 & 0.000 $\pm$ 0.000 \\
SMD-1-1 & PCA-Triage & 0.7 & 0.039 $\pm$ 0.002 & 1.902 $\pm$ 0.455 & 0.006 $\pm$ 0.000 & 757541.0 $\pm$ 0.0 & 3030164.0 $\pm$ 0.0 & 0.700 $\pm$ 0.000 & 0.000 $\pm$ 0.000 \\
SMD-1-1 & Variance & 0.7 & 0.015 $\pm$ 0.000 & 0.852 $\pm$ 0.065 & 0.004 $\pm$ 0.000 & 757541.0 $\pm$ 0.0 & 3030164.0 $\pm$ 0.0 & 0.700 $\pm$ 0.000 & 0.000 $\pm$ 0.000 \\
SMD-1-1 & Threshold & 0.7 & 0.021 $\pm$ 0.000 & 0.562 $\pm$ 0.055 & 0.004 $\pm$ 0.000 & 757541.0 $\pm$ 0.0 & 3030164.0 $\pm$ 0.0 & 0.700 $\pm$ 0.000 & 0.000 $\pm$ 0.000 \\
SMD-1-1 & Uniform & 0.7 & 0.219 $\pm$ 0.043 & 55.388 $\pm$ 21.253 & 0.009 $\pm$ 0.000 & 757541.0 $\pm$ 0.0 & 3030164.0 $\pm$ 0.0 & 0.700 $\pm$ 0.000 & 0.000 $\pm$ 0.000 \\
SMD-1-1 & Full Data & 1.0 & 0.000 $\pm$ 0.000 & 0.000 $\pm$ 0.000 & 0.000 $\pm$ 0.000 & 1082202.0 $\pm$ 0.0 & 4328808.0 $\pm$ 0.0 & 1.000 $\pm$ 0.000 & 0.000 $\pm$ 0.000 \\
\addlinespace
\bottomrule
\end{tabular}%
}
\label{supp-dist-smd-fixed_full_data}
\end{table}
\clearpage
\begin{table}[p]
\centering
\small
\caption{PSM: method retrained. Detection metrics; delay is conditional and measured in samples.}
\resizebox{\textwidth}{!}{%
\begin{tabular}{@{}llccccccccccc@{}}
\toprule
Dataset & Method & Budget & Precision & Recall & FPR & Event recall & Missed events & Delay (samples) & FA/1000 & Point F1 & Specificity & Events \\
\midrule
PSM & PCA-Triage & 0.3 & 0.992 $\pm$ 0.011 & 0.015 $\pm$ 0.001 & 0.000 $\pm$ 0.000 & 0.094 $\pm$ 0.025 & 65.2 $\pm$ 1.8 & 826.403 $\pm$ 131.663 & 0.023 $\pm$ 0.027 & 0.030 $\pm$ 0.002 & 1.000 $\pm$ 0.000 & 72.0 $\pm$ 0.0 \\
PSM & Variance & 0.3 & 0.985 $\pm$ 0.012 & 0.016 $\pm$ 0.001 & 0.000 $\pm$ 0.000 & 0.122 $\pm$ 0.023 & 63.2 $\pm$ 1.6 & 678.380 $\pm$ 299.926 & 0.025 $\pm$ 0.019 & 0.032 $\pm$ 0.002 & 1.000 $\pm$ 0.000 & 72.0 $\pm$ 0.0 \\
PSM & Threshold & 0.3 & 0.974 $\pm$ 0.024 & 0.017 $\pm$ 0.002 & 0.000 $\pm$ 0.000 & 0.103 $\pm$ 0.027 & 64.6 $\pm$ 1.9 & 813.958 $\pm$ 274.543 & 0.046 $\pm$ 0.037 & 0.033 $\pm$ 0.004 & 1.000 $\pm$ 0.000 & 72.0 $\pm$ 0.0 \\
PSM & Uniform & 0.3 & 0.974 $\pm$ 0.030 & 0.016 $\pm$ 0.000 & 0.000 $\pm$ 0.000 & 0.103 $\pm$ 0.021 & 64.6 $\pm$ 1.5 & 683.684 $\pm$ 137.859 & 0.034 $\pm$ 0.032 & 0.031 $\pm$ 0.001 & 1.000 $\pm$ 0.000 & 72.0 $\pm$ 0.0 \\
PSM & PCA-Triage & 0.5 & 0.976 $\pm$ 0.031 & 0.016 $\pm$ 0.001 & 0.000 $\pm$ 0.000 & 0.139 $\pm$ 0.026 & 62.0 $\pm$ 1.9 & 906.807 $\pm$ 359.620 & 0.046 $\pm$ 0.046 & 0.031 $\pm$ 0.001 & 1.000 $\pm$ 0.000 & 72.0 $\pm$ 0.0 \\
PSM & Variance & 0.5 & 0.987 $\pm$ 0.014 & 0.016 $\pm$ 0.001 & 0.000 $\pm$ 0.000 & 0.153 $\pm$ 0.022 & 61.0 $\pm$ 1.6 & 952.049 $\pm$ 214.155 & 0.036 $\pm$ 0.025 & 0.032 $\pm$ 0.001 & 1.000 $\pm$ 0.000 & 72.0 $\pm$ 0.0 \\
PSM & Threshold & 0.5 & 0.980 $\pm$ 0.012 & 0.016 $\pm$ 0.001 & 0.000 $\pm$ 0.000 & 0.125 $\pm$ 0.028 & 63.0 $\pm$ 2.0 & 985.411 $\pm$ 323.244 & 0.064 $\pm$ 0.042 & 0.032 $\pm$ 0.002 & 1.000 $\pm$ 0.000 & 72.0 $\pm$ 0.0 \\
PSM & Uniform & 0.5 & 0.974 $\pm$ 0.022 & 0.017 $\pm$ 0.002 & 0.000 $\pm$ 0.000 & 0.147 $\pm$ 0.032 & 61.4 $\pm$ 2.3 & 678.599 $\pm$ 402.486 & 0.068 $\pm$ 0.055 & 0.033 $\pm$ 0.004 & 1.000 $\pm$ 0.000 & 72.0 $\pm$ 0.0 \\
PSM & PCA-Triage & 0.7 & 0.980 $\pm$ 0.016 & 0.016 $\pm$ 0.001 & 0.000 $\pm$ 0.000 & 0.167 $\pm$ 0.017 & 60.0 $\pm$ 1.2 & 742.202 $\pm$ 405.506 & 0.073 $\pm$ 0.040 & 0.032 $\pm$ 0.001 & 1.000 $\pm$ 0.000 & 72.0 $\pm$ 0.0 \\
PSM & Variance & 0.7 & 0.939 $\pm$ 0.054 & 0.016 $\pm$ 0.001 & 0.000 $\pm$ 0.000 & 0.183 $\pm$ 0.012 & 58.8 $\pm$ 0.8 & 692.545 $\pm$ 110.209 & 0.098 $\pm$ 0.064 & 0.032 $\pm$ 0.002 & 1.000 $\pm$ 0.000 & 72.0 $\pm$ 0.0 \\
PSM & Threshold & 0.7 & 0.980 $\pm$ 0.022 & 0.018 $\pm$ 0.005 & 0.000 $\pm$ 0.000 & 0.181 $\pm$ 0.040 & 59.0 $\pm$ 2.9 & 503.929 $\pm$ 136.598 & 0.109 $\pm$ 0.147 & 0.036 $\pm$ 0.009 & 1.000 $\pm$ 0.000 & 72.0 $\pm$ 0.0 \\
PSM & Uniform & 0.7 & 0.987 $\pm$ 0.013 & 0.016 $\pm$ 0.001 & 0.000 $\pm$ 0.000 & 0.150 $\pm$ 0.033 & 61.2 $\pm$ 2.4 & 676.210 $\pm$ 295.402 & 0.048 $\pm$ 0.047 & 0.032 $\pm$ 0.002 & 1.000 $\pm$ 0.000 & 72.0 $\pm$ 0.0 \\
PSM & Full Data & 1.0 & 0.963 $\pm$ 0.026 & 0.018 $\pm$ 0.004 & 0.000 $\pm$ 0.000 & 0.214 $\pm$ 0.016 & 56.6 $\pm$ 1.1 & 629.185 $\pm$ 95.826 & 0.155 $\pm$ 0.096 & 0.036 $\pm$ 0.007 & 1.000 $\pm$ 0.000 & 72.0 $\pm$ 0.0 \\
\addlinespace
\bottomrule
\end{tabular}%
}
\label{supp-det-psm-method_retrained}
\end{table}
\begin{table}[p]
\centering
\small
\caption{PSM: method retrained. Distortion and mask-counted payload. Four bytes per value are assumed. RMSE uses standardized values.}
\resizebox{\textwidth}{!}{%
\begin{tabular}{@{}llcccccccc@{}}
\toprule
Dataset & Method & Budget & MAE & RMSE & NRMSE & Values & Payload bytes & Realized fraction & Max excess values \\
\midrule
PSM & PCA-Triage & 0.3 & 0.109 $\pm$ 0.000 & 0.290 $\pm$ 0.009 & 0.017 $\pm$ 0.000 & 658807.0 $\pm$ 0.0 & 2635228.0 $\pm$ 0.0 & 0.300 $\pm$ 0.000 & 0.000 $\pm$ 0.000 \\
PSM & Variance & 0.3 & 0.087 $\pm$ 0.000 & 0.255 $\pm$ 0.027 & 0.017 $\pm$ 0.000 & 658807.0 $\pm$ 0.0 & 2635228.0 $\pm$ 0.0 & 0.300 $\pm$ 0.000 & 0.000 $\pm$ 0.000 \\
PSM & Threshold & 0.3 & 0.101 $\pm$ 0.001 & 0.281 $\pm$ 0.009 & 0.019 $\pm$ 0.000 & 658807.0 $\pm$ 0.0 & 2635228.0 $\pm$ 0.0 & 0.300 $\pm$ 0.000 & 0.000 $\pm$ 0.000 \\
PSM & Uniform & 0.3 & 0.105 $\pm$ 0.000 & 0.297 $\pm$ 0.012 & 0.015 $\pm$ 0.000 & 658807.0 $\pm$ 0.0 & 2635228.0 $\pm$ 0.0 & 0.300 $\pm$ 0.000 & 0.000 $\pm$ 0.000 \\
PSM & PCA-Triage & 0.5 & 0.069 $\pm$ 0.000 & 0.222 $\pm$ 0.011 & 0.013 $\pm$ 0.000 & 1098012.0 $\pm$ 0.0 & 4392048.0 $\pm$ 0.0 & 0.500 $\pm$ 0.000 & 0.000 $\pm$ 0.000 \\
PSM & Variance & 0.5 & 0.045 $\pm$ 0.000 & 0.193 $\pm$ 0.011 & 0.013 $\pm$ 0.000 & 1098012.0 $\pm$ 0.0 & 4392048.0 $\pm$ 0.0 & 0.500 $\pm$ 0.000 & 0.000 $\pm$ 0.000 \\
PSM & Threshold & 0.5 & 0.062 $\pm$ 0.000 & 0.228 $\pm$ 0.014 & 0.016 $\pm$ 0.000 & 1098012.0 $\pm$ 0.0 & 4392048.0 $\pm$ 0.0 & 0.500 $\pm$ 0.000 & 0.000 $\pm$ 0.000 \\
PSM & Uniform & 0.5 & 0.069 $\pm$ 0.000 & 0.237 $\pm$ 0.018 & 0.012 $\pm$ 0.000 & 1098012.0 $\pm$ 0.0 & 4392048.0 $\pm$ 0.0 & 0.500 $\pm$ 0.000 & 0.000 $\pm$ 0.000 \\
PSM & PCA-Triage & 0.7 & 0.033 $\pm$ 0.000 & 0.145 $\pm$ 0.010 & 0.010 $\pm$ 0.000 & 1537217.0 $\pm$ 0.0 & 6148868.0 $\pm$ 0.0 & 0.700 $\pm$ 0.000 & 0.000 $\pm$ 0.000 \\
PSM & Variance & 0.7 & 0.015 $\pm$ 0.000 & 0.161 $\pm$ 0.034 & 0.008 $\pm$ 0.000 & 1537217.0 $\pm$ 0.0 & 6148868.0 $\pm$ 0.0 & 0.700 $\pm$ 0.000 & 0.000 $\pm$ 0.000 \\
PSM & Threshold & 0.7 & 0.021 $\pm$ 0.000 & 0.131 $\pm$ 0.014 & 0.008 $\pm$ 0.000 & 1537217.0 $\pm$ 0.0 & 6148868.0 $\pm$ 0.0 & 0.700 $\pm$ 0.000 & 0.000 $\pm$ 0.000 \\
PSM & Uniform & 0.7 & 0.039 $\pm$ 0.000 & 0.176 $\pm$ 0.013 & 0.009 $\pm$ 0.000 & 1537217.0 $\pm$ 0.0 & 6148868.0 $\pm$ 0.0 & 0.700 $\pm$ 0.000 & 0.000 $\pm$ 0.000 \\
PSM & Full Data & 1.0 & 0.000 $\pm$ 0.000 & 0.000 $\pm$ 0.000 & 0.000 $\pm$ 0.000 & 2196025.0 $\pm$ 0.0 & 8784100.0 $\pm$ 0.0 & 1.000 $\pm$ 0.000 & 0.000 $\pm$ 0.000 \\
\addlinespace
\bottomrule
\end{tabular}%
}
\label{supp-dist-psm-method_retrained}
\end{table}
\clearpage
\begin{table}[p]
\centering
\small
\caption{PSM: fixed full data. Detection metrics; delay is conditional and measured in samples.}
\resizebox{\textwidth}{!}{%
\begin{tabular}{@{}llccccccccccc@{}}
\toprule
Dataset & Method & Budget & Precision & Recall & FPR & Event recall & Missed events & Delay (samples) & FA/1000 & Point F1 & Specificity & Events \\
\midrule
PSM & PCA-Triage & 0.3 & 0.973 $\pm$ 0.020 & 0.018 $\pm$ 0.004 & 0.000 $\pm$ 0.000 & 0.122 $\pm$ 0.021 & 63.2 $\pm$ 1.5 & 967.344 $\pm$ 270.657 & 0.050 $\pm$ 0.030 & 0.035 $\pm$ 0.007 & 1.000 $\pm$ 0.000 & 72.0 $\pm$ 0.0 \\
PSM & Variance & 0.3 & 0.982 $\pm$ 0.013 & 0.018 $\pm$ 0.004 & 0.000 $\pm$ 0.000 & 0.150 $\pm$ 0.032 & 61.2 $\pm$ 2.3 & 996.798 $\pm$ 258.641 & 0.039 $\pm$ 0.024 & 0.035 $\pm$ 0.008 & 1.000 $\pm$ 0.000 & 72.0 $\pm$ 0.0 \\
PSM & Threshold & 0.3 & 0.980 $\pm$ 0.015 & 0.018 $\pm$ 0.004 & 0.000 $\pm$ 0.000 & 0.117 $\pm$ 0.025 & 63.6 $\pm$ 1.8 & 1311.094 $\pm$ 265.481 & 0.050 $\pm$ 0.045 & 0.035 $\pm$ 0.008 & 1.000 $\pm$ 0.000 & 72.0 $\pm$ 0.0 \\
PSM & Uniform & 0.3 & 0.970 $\pm$ 0.026 & 0.018 $\pm$ 0.004 & 0.000 $\pm$ 0.000 & 0.122 $\pm$ 0.030 & 63.2 $\pm$ 2.2 & 958.219 $\pm$ 280.034 & 0.055 $\pm$ 0.041 & 0.035 $\pm$ 0.008 & 1.000 $\pm$ 0.000 & 72.0 $\pm$ 0.0 \\
PSM & PCA-Triage & 0.5 & 0.970 $\pm$ 0.023 & 0.018 $\pm$ 0.004 & 0.000 $\pm$ 0.000 & 0.144 $\pm$ 0.021 & 61.6 $\pm$ 1.5 & 1053.611 $\pm$ 302.410 & 0.080 $\pm$ 0.071 & 0.035 $\pm$ 0.007 & 1.000 $\pm$ 0.000 & 72.0 $\pm$ 0.0 \\
PSM & Variance & 0.5 & 0.974 $\pm$ 0.018 & 0.018 $\pm$ 0.004 & 0.000 $\pm$ 0.000 & 0.164 $\pm$ 0.021 & 60.2 $\pm$ 1.5 & 914.783 $\pm$ 154.680 & 0.089 $\pm$ 0.058 & 0.035 $\pm$ 0.007 & 1.000 $\pm$ 0.000 & 72.0 $\pm$ 0.0 \\
PSM & Threshold & 0.5 & 0.972 $\pm$ 0.018 & 0.018 $\pm$ 0.003 & 0.000 $\pm$ 0.000 & 0.136 $\pm$ 0.050 & 62.2 $\pm$ 3.6 & 852.212 $\pm$ 339.140 & 0.096 $\pm$ 0.067 & 0.034 $\pm$ 0.007 & 1.000 $\pm$ 0.000 & 72.0 $\pm$ 0.0 \\
PSM & Uniform & 0.5 & 0.971 $\pm$ 0.028 & 0.018 $\pm$ 0.004 & 0.000 $\pm$ 0.000 & 0.150 $\pm$ 0.023 & 61.2 $\pm$ 1.6 & 789.026 $\pm$ 199.579 & 0.061 $\pm$ 0.061 & 0.035 $\pm$ 0.007 & 1.000 $\pm$ 0.000 & 72.0 $\pm$ 0.0 \\
PSM & PCA-Triage & 0.7 & 0.969 $\pm$ 0.024 & 0.018 $\pm$ 0.003 & 0.000 $\pm$ 0.000 & 0.181 $\pm$ 0.026 & 59.0 $\pm$ 1.9 & 758.742 $\pm$ 54.960 & 0.112 $\pm$ 0.065 & 0.035 $\pm$ 0.007 & 1.000 $\pm$ 0.000 & 72.0 $\pm$ 0.0 \\
PSM & Variance & 0.7 & 0.969 $\pm$ 0.024 & 0.018 $\pm$ 0.004 & 0.000 $\pm$ 0.000 & 0.183 $\pm$ 0.021 & 58.8 $\pm$ 1.5 & 723.110 $\pm$ 142.210 & 0.116 $\pm$ 0.079 & 0.035 $\pm$ 0.007 & 1.000 $\pm$ 0.000 & 72.0 $\pm$ 0.0 \\
PSM & Threshold & 0.7 & 0.967 $\pm$ 0.027 & 0.018 $\pm$ 0.003 & 0.000 $\pm$ 0.000 & 0.192 $\pm$ 0.021 & 58.2 $\pm$ 1.5 & 719.205 $\pm$ 119.058 & 0.121 $\pm$ 0.080 & 0.036 $\pm$ 0.006 & 1.000 $\pm$ 0.000 & 72.0 $\pm$ 0.0 \\
PSM & Uniform & 0.7 & 0.974 $\pm$ 0.026 & 0.018 $\pm$ 0.003 & 0.000 $\pm$ 0.000 & 0.169 $\pm$ 0.015 & 59.8 $\pm$ 1.1 & 762.339 $\pm$ 170.357 & 0.093 $\pm$ 0.072 & 0.035 $\pm$ 0.006 & 1.000 $\pm$ 0.000 & 72.0 $\pm$ 0.0 \\
PSM & Full Data & 1.0 & 0.963 $\pm$ 0.026 & 0.018 $\pm$ 0.004 & 0.000 $\pm$ 0.000 & 0.214 $\pm$ 0.016 & 56.6 $\pm$ 1.1 & 629.185 $\pm$ 95.826 & 0.155 $\pm$ 0.096 & 0.036 $\pm$ 0.007 & 1.000 $\pm$ 0.000 & 72.0 $\pm$ 0.0 \\
\addlinespace
\bottomrule
\end{tabular}%
}
\label{supp-det-psm-fixed_full_data}
\end{table}
\begin{table}[p]
\centering
\small
\caption{PSM: fixed full data. Distortion and mask-counted payload. Four bytes per value are assumed. RMSE uses standardized values.}
\resizebox{\textwidth}{!}{%
\begin{tabular}{@{}llcccccccc@{}}
\toprule
Dataset & Method & Budget & MAE & RMSE & NRMSE & Values & Payload bytes & Realized fraction & Max excess values \\
\midrule
PSM & PCA-Triage & 0.3 & 0.109 $\pm$ 0.000 & 0.290 $\pm$ 0.009 & 0.017 $\pm$ 0.000 & 658807.0 $\pm$ 0.0 & 2635228.0 $\pm$ 0.0 & 0.300 $\pm$ 0.000 & 0.000 $\pm$ 0.000 \\
PSM & Variance & 0.3 & 0.087 $\pm$ 0.000 & 0.255 $\pm$ 0.027 & 0.017 $\pm$ 0.000 & 658807.0 $\pm$ 0.0 & 2635228.0 $\pm$ 0.0 & 0.300 $\pm$ 0.000 & 0.000 $\pm$ 0.000 \\
PSM & Threshold & 0.3 & 0.101 $\pm$ 0.001 & 0.281 $\pm$ 0.009 & 0.019 $\pm$ 0.000 & 658807.0 $\pm$ 0.0 & 2635228.0 $\pm$ 0.0 & 0.300 $\pm$ 0.000 & 0.000 $\pm$ 0.000 \\
PSM & Uniform & 0.3 & 0.105 $\pm$ 0.000 & 0.297 $\pm$ 0.012 & 0.015 $\pm$ 0.000 & 658807.0 $\pm$ 0.0 & 2635228.0 $\pm$ 0.0 & 0.300 $\pm$ 0.000 & 0.000 $\pm$ 0.000 \\
PSM & PCA-Triage & 0.5 & 0.069 $\pm$ 0.000 & 0.222 $\pm$ 0.011 & 0.013 $\pm$ 0.000 & 1098012.0 $\pm$ 0.0 & 4392048.0 $\pm$ 0.0 & 0.500 $\pm$ 0.000 & 0.000 $\pm$ 0.000 \\
PSM & Variance & 0.5 & 0.045 $\pm$ 0.000 & 0.193 $\pm$ 0.011 & 0.013 $\pm$ 0.000 & 1098012.0 $\pm$ 0.0 & 4392048.0 $\pm$ 0.0 & 0.500 $\pm$ 0.000 & 0.000 $\pm$ 0.000 \\
PSM & Threshold & 0.5 & 0.062 $\pm$ 0.000 & 0.228 $\pm$ 0.014 & 0.016 $\pm$ 0.000 & 1098012.0 $\pm$ 0.0 & 4392048.0 $\pm$ 0.0 & 0.500 $\pm$ 0.000 & 0.000 $\pm$ 0.000 \\
PSM & Uniform & 0.5 & 0.069 $\pm$ 0.000 & 0.237 $\pm$ 0.018 & 0.012 $\pm$ 0.000 & 1098012.0 $\pm$ 0.0 & 4392048.0 $\pm$ 0.0 & 0.500 $\pm$ 0.000 & 0.000 $\pm$ 0.000 \\
PSM & PCA-Triage & 0.7 & 0.033 $\pm$ 0.000 & 0.145 $\pm$ 0.010 & 0.010 $\pm$ 0.000 & 1537217.0 $\pm$ 0.0 & 6148868.0 $\pm$ 0.0 & 0.700 $\pm$ 0.000 & 0.000 $\pm$ 0.000 \\
PSM & Variance & 0.7 & 0.015 $\pm$ 0.000 & 0.161 $\pm$ 0.034 & 0.008 $\pm$ 0.000 & 1537217.0 $\pm$ 0.0 & 6148868.0 $\pm$ 0.0 & 0.700 $\pm$ 0.000 & 0.000 $\pm$ 0.000 \\
PSM & Threshold & 0.7 & 0.021 $\pm$ 0.000 & 0.131 $\pm$ 0.014 & 0.008 $\pm$ 0.000 & 1537217.0 $\pm$ 0.0 & 6148868.0 $\pm$ 0.0 & 0.700 $\pm$ 0.000 & 0.000 $\pm$ 0.000 \\
PSM & Uniform & 0.7 & 0.039 $\pm$ 0.000 & 0.176 $\pm$ 0.013 & 0.009 $\pm$ 0.000 & 1537217.0 $\pm$ 0.0 & 6148868.0 $\pm$ 0.0 & 0.700 $\pm$ 0.000 & 0.000 $\pm$ 0.000 \\
PSM & Full Data & 1.0 & 0.000 $\pm$ 0.000 & 0.000 $\pm$ 0.000 & 0.000 $\pm$ 0.000 & 2196025.0 $\pm$ 0.0 & 8784100.0 $\pm$ 0.0 & 1.000 $\pm$ 0.000 & 0.000 $\pm$ 0.000 \\
\addlinespace
\bottomrule
\end{tabular}%
}
\label{supp-dist-psm-fixed_full_data}
\end{table}
\clearpage
\begin{table}[p]
\centering
\small
\caption{TEP: method retrained. Detection metrics; delay is conditional and measured in samples.}
\resizebox{\textwidth}{!}{%
\begin{tabular}{@{}llccccccccccc@{}}
\toprule
Dataset & Method & Budget & Precision & Recall & FPR & Event recall & Missed events & Delay (samples) & FA/1000 & Point F1 & Specificity & Events \\
\midrule
TEP & PCA-Triage & 0.5 & 0.963 $\pm$ 0.001 & 0.904 $\pm$ 0.005 & 0.384 $\pm$ 0.012 & 1.000 $\pm$ 0.000 & 0.0 $\pm$ 0.0 & 2.962 $\pm$ 0.879 & 7.223 $\pm$ 0.391 & 0.932 $\pm$ 0.002 & 0.616 $\pm$ 0.012 & 100.0 $\pm$ 0.0 \\
TEP & Variance & 0.5 & 0.962 $\pm$ 0.001 & 0.903 $\pm$ 0.004 & 0.393 $\pm$ 0.009 & 1.000 $\pm$ 0.000 & 0.0 $\pm$ 0.0 & 3.460 $\pm$ 0.554 & 7.170 $\pm$ 0.668 & 0.931 $\pm$ 0.002 & 0.607 $\pm$ 0.009 & 100.0 $\pm$ 0.0 \\
TEP & Threshold & 0.5 & 0.965 $\pm$ 0.002 & 0.886 $\pm$ 0.007 & 0.354 $\pm$ 0.023 & 1.000 $\pm$ 0.000 & 0.0 $\pm$ 0.0 & 4.588 $\pm$ 0.696 & 5.760 $\pm$ 0.398 & 0.924 $\pm$ 0.003 & 0.646 $\pm$ 0.023 & 100.0 $\pm$ 0.0 \\
TEP & Uniform & 0.5 & 0.962 $\pm$ 0.001 & 0.914 $\pm$ 0.002 & 0.399 $\pm$ 0.011 & 1.000 $\pm$ 0.000 & 0.0 $\pm$ 0.0 & 2.044 $\pm$ 0.438 & 9.672 $\pm$ 0.148 & 0.937 $\pm$ 0.001 & 0.601 $\pm$ 0.011 & 100.0 $\pm$ 0.0 \\
TEP & Full Data & 1.0 & 0.943 $\pm$ 0.002 & 0.972 $\pm$ 0.002 & 0.638 $\pm$ 0.019 & 1.000 $\pm$ 0.000 & 0.0 $\pm$ 0.0 & 0.350 $\pm$ 0.304 & 8.343 $\pm$ 1.114 & 0.958 $\pm$ 0.001 & 0.362 $\pm$ 0.019 & 100.0 $\pm$ 0.0 \\
\addlinespace
\bottomrule
\end{tabular}%
}
\label{supp-det-tep-method_retrained}
\end{table}
\begin{table}[p]
\centering
\small
\caption{TEP: method retrained. Distortion and mask-counted payload. Four bytes per value are assumed. RMSE uses standardized values.}
\resizebox{\textwidth}{!}{%
\begin{tabular}{@{}llcccccccc@{}}
\toprule
Dataset & Method & Budget & MAE & RMSE & NRMSE & Values & Payload bytes & Realized fraction & Max excess values \\
\midrule
TEP & PCA-Triage & 0.5 & 0.104 $\pm$ 0.000 & 0.270 $\pm$ 0.002 & 0.027 $\pm$ 0.000 & 1365000.0 $\pm$ 0.0 & 5460000.0 $\pm$ 0.0 & 0.500 $\pm$ 0.000 & 0.000 $\pm$ 0.000 \\
TEP & Variance & 0.5 & 0.099 $\pm$ 0.000 & 0.261 $\pm$ 0.002 & 0.026 $\pm$ 0.000 & 1365000.0 $\pm$ 0.0 & 5460000.0 $\pm$ 0.0 & 0.500 $\pm$ 0.000 & 0.000 $\pm$ 0.000 \\
TEP & Threshold & 0.5 & 0.139 $\pm$ 0.001 & 0.383 $\pm$ 0.002 & 0.038 $\pm$ 0.000 & 1365000.0 $\pm$ 0.0 & 5460000.0 $\pm$ 0.0 & 0.500 $\pm$ 0.000 & 0.000 $\pm$ 0.000 \\
TEP & Uniform & 0.5 & 0.175 $\pm$ 0.000 & 0.483 $\pm$ 0.001 & 0.036 $\pm$ 0.000 & 1365000.0 $\pm$ 0.0 & 5460000.0 $\pm$ 0.0 & 0.500 $\pm$ 0.000 & 0.000 $\pm$ 0.000 \\
TEP & Full Data & 1.0 & 0.000 $\pm$ 0.000 & 0.000 $\pm$ 0.000 & 0.000 $\pm$ 0.000 & 2730000.0 $\pm$ 0.0 & 10920000.0 $\pm$ 0.0 & 1.000 $\pm$ 0.000 & 0.000 $\pm$ 0.000 \\
\addlinespace
\bottomrule
\end{tabular}%
}
\label{supp-dist-tep-method_retrained}
\end{table}
\clearpage
\begin{table}[p]
\centering
\small
\caption{TEP: fixed full data. Detection metrics; delay is conditional and measured in samples.}
\resizebox{\textwidth}{!}{%
\begin{tabular}{@{}llccccccccccc@{}}
\toprule
Dataset & Method & Budget & Precision & Recall & FPR & Event recall & Missed events & Delay (samples) & FA/1000 & Point F1 & Specificity & Events \\
\midrule
TEP & PCA-Triage & 0.5 & 0.940 $\pm$ 0.001 & 0.980 $\pm$ 0.001 & 0.688 $\pm$ 0.014 & 1.000 $\pm$ 0.000 & 0.0 $\pm$ 0.0 & 0.242 $\pm$ 0.094 & 7.878 $\pm$ 0.378 & 0.960 $\pm$ 0.001 & 0.312 $\pm$ 0.014 & 100.0 $\pm$ 0.0 \\
TEP & Variance & 0.5 & 0.939 $\pm$ 0.001 & 0.980 $\pm$ 0.001 & 0.691 $\pm$ 0.017 & 1.000 $\pm$ 0.000 & 0.0 $\pm$ 0.0 & 0.242 $\pm$ 0.094 & 7.779 $\pm$ 0.274 & 0.959 $\pm$ 0.000 & 0.309 $\pm$ 0.017 & 100.0 $\pm$ 0.0 \\
TEP & Threshold & 0.5 & 0.939 $\pm$ 0.001 & 0.983 $\pm$ 0.001 & 0.698 $\pm$ 0.012 & 1.000 $\pm$ 0.000 & 0.0 $\pm$ 0.0 & 0.242 $\pm$ 0.094 & 7.333 $\pm$ 0.146 & 0.960 $\pm$ 0.001 & 0.302 $\pm$ 0.012 & 100.0 $\pm$ 0.0 \\
TEP & Uniform & 0.5 & 0.940 $\pm$ 0.001 & 0.975 $\pm$ 0.001 & 0.678 $\pm$ 0.015 & 1.000 $\pm$ 0.000 & 0.0 $\pm$ 0.0 & 0.242 $\pm$ 0.094 & 8.560 $\pm$ 0.149 & 0.958 $\pm$ 0.001 & 0.322 $\pm$ 0.015 & 100.0 $\pm$ 0.0 \\
TEP & Full Data & 1.0 & 0.943 $\pm$ 0.002 & 0.972 $\pm$ 0.002 & 0.638 $\pm$ 0.019 & 1.000 $\pm$ 0.000 & 0.0 $\pm$ 0.0 & 0.350 $\pm$ 0.304 & 8.343 $\pm$ 1.114 & 0.958 $\pm$ 0.001 & 0.362 $\pm$ 0.019 & 100.0 $\pm$ 0.0 \\
\addlinespace
\bottomrule
\end{tabular}%
}
\label{supp-det-tep-fixed_full_data}
\end{table}
\begin{table}[p]
\centering
\small
\caption{TEP: fixed full data. Distortion and mask-counted payload. Four bytes per value are assumed. RMSE uses standardized values.}
\resizebox{\textwidth}{!}{%
\begin{tabular}{@{}llcccccccc@{}}
\toprule
Dataset & Method & Budget & MAE & RMSE & NRMSE & Values & Payload bytes & Realized fraction & Max excess values \\
\midrule
TEP & PCA-Triage & 0.5 & 0.104 $\pm$ 0.000 & 0.270 $\pm$ 0.002 & 0.027 $\pm$ 0.000 & 1365000.0 $\pm$ 0.0 & 5460000.0 $\pm$ 0.0 & 0.500 $\pm$ 0.000 & 0.000 $\pm$ 0.000 \\
TEP & Variance & 0.5 & 0.099 $\pm$ 0.000 & 0.261 $\pm$ 0.002 & 0.026 $\pm$ 0.000 & 1365000.0 $\pm$ 0.0 & 5460000.0 $\pm$ 0.0 & 0.500 $\pm$ 0.000 & 0.000 $\pm$ 0.000 \\
TEP & Threshold & 0.5 & 0.139 $\pm$ 0.001 & 0.383 $\pm$ 0.002 & 0.038 $\pm$ 0.000 & 1365000.0 $\pm$ 0.0 & 5460000.0 $\pm$ 0.0 & 0.500 $\pm$ 0.000 & 0.000 $\pm$ 0.000 \\
TEP & Uniform & 0.5 & 0.175 $\pm$ 0.000 & 0.483 $\pm$ 0.001 & 0.036 $\pm$ 0.000 & 1365000.0 $\pm$ 0.0 & 5460000.0 $\pm$ 0.0 & 0.500 $\pm$ 0.000 & 0.000 $\pm$ 0.000 \\
TEP & Full Data & 1.0 & 0.000 $\pm$ 0.000 & 0.000 $\pm$ 0.000 & 0.000 $\pm$ 0.000 & 2730000.0 $\pm$ 0.0 & 10920000.0 $\pm$ 0.0 & 1.000 $\pm$ 0.000 & 0.000 $\pm$ 0.000 \\
\addlinespace
\bottomrule
\end{tabular}%
}
\label{supp-dist-tep-fixed_full_data}
\end{table}
\clearpage

\begin{table}[ht]
\centering
\small
\caption{SMD: method retrained. Point F1 at 50\% bandwidth across predeclared training-score quantiles; Full Data is the full-rate reference. The primary quantile remains 0.99.}
\begin{tabular}{lccccc}
\toprule
Method & 0.95 & 0.975 & 0.99 & 0.995 & 0.999 \\
\midrule
PCA-Triage & 0.389 $\pm$ 0.009 & 0.423 $\pm$ 0.017 & 0.390 $\pm$ 0.017 & 0.369 $\pm$ 0.004 & 0.355 $\pm$ 0.003 \\
Variance & 0.375 $\pm$ 0.022 & 0.417 $\pm$ 0.019 & 0.407 $\pm$ 0.019 & 0.390 $\pm$ 0.016 & 0.357 $\pm$ 0.004 \\
Threshold & 0.392 $\pm$ 0.023 & 0.427 $\pm$ 0.018 & 0.393 $\pm$ 0.013 & 0.371 $\pm$ 0.008 & 0.350 $\pm$ 0.001 \\
Uniform & 0.377 $\pm$ 0.029 & 0.430 $\pm$ 0.014 & 0.395 $\pm$ 0.008 & 0.373 $\pm$ 0.005 & 0.357 $\pm$ 0.002 \\
Full Data & 0.383 $\pm$ 0.018 & 0.431 $\pm$ 0.027 & 0.409 $\pm$ 0.010 & 0.388 $\pm$ 0.009 & 0.356 $\pm$ 0.002 \\
\bottomrule
\end{tabular}
\end{table}
\begin{table}[ht]
\centering
\small
\caption{SMD: fixed full data. Point F1 at 50\% bandwidth across predeclared training-score quantiles; Full Data is the full-rate reference. The primary quantile remains 0.99.}
\begin{tabular}{lccccc}
\toprule
Method & 0.95 & 0.975 & 0.99 & 0.995 & 0.999 \\
\midrule
PCA-Triage & 0.386 $\pm$ 0.019 & 0.448 $\pm$ 0.026 & 0.401 $\pm$ 0.012 & 0.370 $\pm$ 0.002 & 0.357 $\pm$ 0.001 \\
Variance & 0.382 $\pm$ 0.019 & 0.433 $\pm$ 0.026 & 0.406 $\pm$ 0.011 & 0.385 $\pm$ 0.009 & 0.356 $\pm$ 0.002 \\
Threshold & 0.374 $\pm$ 0.018 & 0.429 $\pm$ 0.027 & 0.396 $\pm$ 0.013 & 0.375 $\pm$ 0.010 & 0.349 $\pm$ 0.003 \\
Uniform & 0.385 $\pm$ 0.020 & 0.446 $\pm$ 0.027 & 0.400 $\pm$ 0.010 & 0.375 $\pm$ 0.006 & 0.356 $\pm$ 0.002 \\
Full Data & 0.383 $\pm$ 0.018 & 0.431 $\pm$ 0.027 & 0.409 $\pm$ 0.010 & 0.388 $\pm$ 0.009 & 0.356 $\pm$ 0.002 \\
\bottomrule
\end{tabular}
\end{table}
\clearpage
\begin{table}[ht]
\centering
\small
\caption{PSM: method retrained. Point F1 at 50\% bandwidth across predeclared training-score quantiles; Full Data is the full-rate reference. The primary quantile remains 0.99.}
\begin{tabular}{lccccc}
\toprule
Method & 0.95 & 0.975 & 0.99 & 0.995 & 0.999 \\
\midrule
PCA-Triage & 0.277 $\pm$ 0.047 & 0.104 $\pm$ 0.011 & 0.031 $\pm$ 0.001 & 0.027 $\pm$ 0.000 & 0.023 $\pm$ 0.000 \\
Variance & 0.267 $\pm$ 0.072 & 0.098 $\pm$ 0.032 & 0.032 $\pm$ 0.001 & 0.027 $\pm$ 0.000 & 0.023 $\pm$ 0.000 \\
Threshold & 0.237 $\pm$ 0.052 & 0.088 $\pm$ 0.016 & 0.032 $\pm$ 0.002 & 0.027 $\pm$ 0.000 & 0.024 $\pm$ 0.001 \\
Uniform & 0.268 $\pm$ 0.061 & 0.102 $\pm$ 0.029 & 0.033 $\pm$ 0.004 & 0.027 $\pm$ 0.000 & 0.023 $\pm$ 0.001 \\
Full Data & 0.272 $\pm$ 0.047 & 0.108 $\pm$ 0.023 & 0.036 $\pm$ 0.007 & 0.027 $\pm$ 0.000 & 0.024 $\pm$ 0.000 \\
\bottomrule
\end{tabular}
\end{table}
\begin{table}[ht]
\centering
\small
\caption{PSM: fixed full data. Point F1 at 50\% bandwidth across predeclared training-score quantiles; Full Data is the full-rate reference. The primary quantile remains 0.99.}
\begin{tabular}{lccccc}
\toprule
Method & 0.95 & 0.975 & 0.99 & 0.995 & 0.999 \\
\midrule
PCA-Triage & 0.269 $\pm$ 0.049 & 0.105 $\pm$ 0.021 & 0.035 $\pm$ 0.007 & 0.027 $\pm$ 0.000 & 0.023 $\pm$ 0.000 \\
Variance & 0.272 $\pm$ 0.049 & 0.105 $\pm$ 0.023 & 0.035 $\pm$ 0.007 & 0.027 $\pm$ 0.000 & 0.023 $\pm$ 0.000 \\
Threshold & 0.271 $\pm$ 0.052 & 0.104 $\pm$ 0.023 & 0.034 $\pm$ 0.007 & 0.027 $\pm$ 0.000 & 0.022 $\pm$ 0.001 \\
Uniform & 0.270 $\pm$ 0.048 & 0.105 $\pm$ 0.022 & 0.035 $\pm$ 0.007 & 0.027 $\pm$ 0.000 & 0.023 $\pm$ 0.000 \\
Full Data & 0.272 $\pm$ 0.047 & 0.108 $\pm$ 0.023 & 0.036 $\pm$ 0.007 & 0.027 $\pm$ 0.000 & 0.024 $\pm$ 0.000 \\
\bottomrule
\end{tabular}
\end{table}
\clearpage

\end{document}
